\documentclass[11pt]{article}
\usepackage[margin=1in]{geometry}
\usepackage{enumitem}
% =============================================================================
% Preambles/header.tex — shared LaTeX/Beamer preamble for this project
%
% Use from a Beamer lecture by prepending:
%     \input{header}
% and compiling with TEXINPUTS=../Preambles:$TEXINPUTS (the /compile-latex
% skill does this automatically).
%
% PALETTE CONTRACT. Color names defined here MUST match the names in
% Quarto/theme-template.scss so Beamer and Quarto renderings use the same
% palette. The scripts/check-palette-sync.sh script enforces this.
%
% To customize: edit the HEX values below AND the matching $variable in
% Quarto/theme-template.scss, then run scripts/check-palette-sync.sh.
% =============================================================================

% -----------------------------------------------------------------------------
% Engine detection + encoding
%
% Our /compile-latex skill uses XeLaTeX, where inputenc is unnecessary and
% can produce encoding warnings. Only load inputenc under pdfTeX.
% -----------------------------------------------------------------------------
\usepackage{iftex}
\ifPDFTeX
  \usepackage[utf8]{inputenc}
\fi

\usepackage{amsmath, amssymb, amsthm}
\usepackage{booktabs}
\usepackage{graphicx}

% -----------------------------------------------------------------------------
% PALETTE — mirrors Quarto/theme-template.scss variable names.
% Load xcolor and define colors BEFORE anything that references them
% (notably hyperref, hypersetup, and the Beamer color assignments).
% -----------------------------------------------------------------------------
\usepackage{xcolor}

\definecolor{primary-blue}{HTML}{012169}      % headings, accents
\definecolor{primary-gold}{HTML}{B9975B}      % emphasis, borders
\definecolor{highlight-yellow}{HTML}{F2A900}  % markers, alerts
\definecolor{light-bg}{HTML}{E8EDF5}          % title / section backgrounds
\definecolor{jet}{HTML}{1A1A1A}               % body text

% Semantic colors (used in diagrams and annotations)
\definecolor{positive}{HTML}{15803D}          % good / observed / identified
\definecolor{negative}{HTML}{B91C1C}          % bad / problematic / bias
\definecolor{neutral}{HTML}{525252}           % reference / context

% Highlight accents (match .hi-* classes in the SCSS)
\definecolor{hi-slate}{HTML}{314F4F}
\definecolor{hi-green}{HTML}{2E8B57}
\definecolor{hi-red}{HTML}{C41E3A}

% Semantic aliases so skill templates and snippets can write
% \textcolor{accent}{...} without hard-coding "primary-blue".
\colorlet{accent}{primary-blue}
\colorlet{accent2}{primary-gold}

% -----------------------------------------------------------------------------
% Hyperref — loaded AFTER xcolor + palette so link colors resolve.
% -----------------------------------------------------------------------------
\usepackage{hyperref}
\hypersetup{colorlinks=true, linkcolor=primary-blue, urlcolor=primary-blue,
            citecolor=primary-blue, pdfborder={0 0 0}}

% -----------------------------------------------------------------------------
% Course code — set once per deck (after \input{header}, before \begin{document})
% via \coursecode{CS401}. Shown in the footer on every slide so a deck's course
% is identifiable without opening the title page. Empty by default.
% -----------------------------------------------------------------------------
\makeatletter
\newcommand{\coursecode}[1]{\gdef\@coursecodetext{#1}}
\gdef\@coursecodetext{}
\makeatother

% -----------------------------------------------------------------------------
% Beamer theme assignments (only apply under Beamer).
%
% We detect Beamer via \@ifclassloaded{beamer}, which is the documented,
% non-internal way. This must be wrapped in \makeatletter/\makeatother
% because \@ifclassloaded contains an @-catcode character.
% -----------------------------------------------------------------------------
\makeatletter
\@ifclassloaded{beamer}{%
  \usefonttheme{professionalfonts}
  \setbeamercolor{structure}{fg=primary-blue}
  \setbeamercolor{frametitle}{fg=primary-blue}
  \setbeamercolor{title}{fg=primary-blue}
  \setbeamercolor{subtitle}{fg=primary-gold}
  \setbeamercolor{author}{fg=primary-blue}
  \setbeamercolor{institute}{fg=neutral}
  \setbeamercolor{date}{fg=primary-gold}
  \setbeamercolor{itemize item}{fg=primary-blue}
  \setbeamercolor{itemize subitem}{fg=primary-gold}
  \setbeamercolor{alerted text}{fg=primary-gold}
  \setbeamercolor{block title}{fg=white, bg=primary-blue}
  \setbeamercolor{block body}{bg=light-bg}
  % Semantic block variants: block=definition/concept (blue), exampleblock=
  % worked example (green/positive), alertblock=key takeaway or caution (gold).
  % Keep to <=2 colored boxes per slide across ALL three variants (INV-7).
  \setbeamercolor{block title example}{fg=white, bg=positive}
  \setbeamercolor{block body example}{bg=light-bg}
  \setbeamercolor{block title alerted}{fg=white, bg=primary-gold}
  \setbeamercolor{block body alerted}{bg=light-bg}

  % Minimal footer: course code (left) + page X / N (right), no navigation clutter
  \setbeamertemplate{navigation symbols}{}
  \setbeamertemplate{footline}{%
    \color{neutral}\footnotesize\hspace{0.7em}\@coursecodetext\hfill
    \insertframenumber/\inserttotalframenumber\hspace{0.7em}\vspace{0.3em}}
}{}
\makeatother

% -----------------------------------------------------------------------------
% TikZ setup — libraries the snippet gallery uses
% -----------------------------------------------------------------------------
\usepackage{tikz}
\usetikzlibrary{arrows.meta, positioning, calc, decorations.pathreplacing,
                fit, shapes.geometric, backgrounds}

% Shared TikZ styles — snippets and hand-written diagrams can pick these up
% via `\begin{tikzpicture}[every node/.style={...}]` or by naming specific
% styles (e.g., `\node[dag-node] (X) at (0,0) {X};`).
\tikzset{
  % Node styles for causal DAGs and similar
  dag-node/.style = {
    draw=primary-blue, fill=light-bg, rounded corners=2pt,
    minimum width=1.2cm, minimum height=0.9cm, align=center, font=\small
  },
  decision-node/.style = {
    draw=primary-gold, fill=white, diamond, aspect=1.8,
    minimum width=1.8cm, minimum height=1.2cm, align=center, font=\small,
    inner sep=1pt
  },
  % Edge styles — observed vs counterfactual
  observed-edge/.style = {-{Latex[length=2.5mm]}, thick, draw=jet},
  counterfactual-edge/.style = {-{Latex[length=2.5mm]}, thick, draw=neutral, dashed},
  confound-edge/.style = {-{Latex[length=2.5mm]}, thick, draw=negative, dashed},
  % Data-point styles for plots
  observed-dot/.style = {fill=primary-blue, draw=primary-blue, circle, inner sep=1.2pt},
  counterfactual-dot/.style = {fill=white, draw=neutral, circle, inner sep=1.2pt}
}

% -----------------------------------------------------------------------------
% Convenience macros
% -----------------------------------------------------------------------------
\newcommand{\muted}[1]{\textcolor{neutral}{#1}}
\newcommand{\key}[1]{\textcolor{primary-gold}{\textbf{#1}}}
\newcommand{\good}[1]{\textcolor{positive}{#1}}
\newcommand{\bad}[1]{\textcolor{negative}{#1}}

% Standout frame macro: full-bleed dark background for section transitions.
% Beamer-only (uses \begin{frame}/\setbeamercolor/\usebeamerfont) -- do not
% call from an article-class file (e.g. Notes/<CODE>/*-notes.tex). \muted,
% \key, \good, \bad, and \coursecode above are plain xcolor/gdef and are
% safe to use from either class; verified when Notes/article-class support
% was added.
% Expand: \transitionslide{Your transition title}
\newcommand{\transitionslide}[1]{%
  \begingroup
  \setbeamercolor{background canvas}{bg=primary-blue}
  \begin{frame}[plain]
    \centering
    \vfill
    {\usebeamerfont{title}\color{white}\Large #1\par}
    \vfill
  \end{frame}
  \endgroup
}

% Section divider frame (Beamer-only), an alternative to \transitionslide with a
% darker two-line style: near-black (jet) background, a small white label line
% above a larger gold title, and a thin gold rule along the bottom edge.
% Expand: \sectiondivider{Section 2}{Pointers}
\newcommand{\sectiondivider}[2]{%
  \begingroup
  \setbeamercolor{background canvas}{bg=jet}
  \begin{frame}[plain]
    \vspace*{3.0cm}
    {\color{white}\ttfamily\large #1\par}
    \vspace{0.5cm}
    {\color{primary-gold}\LARGE #2\par}
    \begin{tikzpicture}[remember picture, overlay]
      \draw[primary-gold, line width=2.5pt]
        ($(current page.south west)+(0,0.1)$) -- ($(current page.south east)+(0,0.1)$);
    \end{tikzpicture}
  \end{frame}
  \endgroup
}

% -----------------------------------------------------------------------------
% Channel watermark — "Concepts That Click" (YouTube).
%
% Article-class files (CompetitiveExam Q/A sets, Notes, Assignments) get a
% light diagonal draft watermark on every page plus a centered footer naming
% the channel. Beamer decks get a subtle TikZ background watermark instead
% (the footline already carries the course code). Centralized here so every
% MCQ PDF is branded without per-file edits.
% -----------------------------------------------------------------------------
\makeatletter
\@ifclassloaded{article}{%
  \usepackage{draftwatermark}
  \SetWatermarkText{Concepts That Click}
  \SetWatermarkScale{1}
  \SetWatermarkFontSize{2cm}
  \SetWatermarkLightness{0.86}
  \SetWatermarkAngle{45}
  \usepackage{fancyhdr}
  \pagestyle{fancy}
  \fancyhf{}
  \fancyfoot[C]{\footnotesize\textcolor{neutral}{Concepts That Click~|~YouTube: \href{https://www.youtube.com/@concepts.thatclick}{@concepts.thatclick}}}
  \renewcommand{\headrulewidth}{0pt}
  \renewcommand{\footrulewidth}{0pt}
  \fancypagestyle{plain}{%
    \fancyhf{}%
    \fancyfoot[C]{\footnotesize\textcolor{neutral}{Concepts That Click~|~YouTube: \href{https://www.youtube.com/@concepts.thatclick}{@concepts.thatclick}}}%
    \renewcommand{\headrulewidth}{0pt}%
    \renewcommand{\footrulewidth}{0pt}%
  }
}{}
\@ifclassloaded{beamer}{%
  \setbeamertemplate{background}{%
    \begin{tikzpicture}[remember picture, overlay]
      \node[rotate=45, opacity=0.055, align=center]
        at (current page.center)
        {\Huge\textcolor{neutral}{Concepts That Click}};
    \end{tikzpicture}%
  }
}{}
\makeatother 
\coursecode{JKSSB}
\renewcommand{\thesection}{46.\arabic{section}}
\newtheorem{example}{Example}[section]
\newenvironment{definitionbox}[1]{%
  \par\noindent\textbf{Definition (#1).}\ \itshape
}{\par\vspace{0.3em}}
\title{Firmware, BIOS, UEFI and Drivers --- Detailed Lecture Notes}
\author{JKSSB Computer Awareness}
\date{\today}

\begin{document}
\maketitle

\section{What firmware is}
Every computer has two broad parts: the hardware, which is the physical
circuitry, and the software, which is the set of instructions. Between them
sits a special kind of software called \textbf{firmware}. Firmware is
software that is stored in non-volatile memory on the hardware itself and
that gives the device its basic, low-level instructions. It is not opened by
the user like a word processor or a web browser; it is fixed in the device
and runs by itself.

Firmware runs \emph{before} the operating system. When a computer is switched
on, its memory is empty and no program is running, so something must tell the
hardware how to begin. That ``something'' is the firmware. It starts the
hardware, checks it, and then hands control to the operating system.

\begin{definitionbox}{Firmware}
Software stored in non-volatile memory that is programmed into a hardware
device and provides its basic, low-level startup and control instructions,
running before the operating system.
\end{definitionbox}

Firmware is different from ordinary application software in three ways. It is
stored on the device rather than installed like a program; it is tied closely
to the exact hardware; and it starts the machine rather than being started by
the user.

\section{Where firmware is stored}
Firmware must be available the instant the machine is switched on, so it is
kept in \textbf{non-volatile} memory --- memory that keeps its contents
without power. If firmware were stored in RAM, it would be lost the moment
the power was switched off, and the machine could never start.

Common homes for firmware are \textbf{ROM} (Read-Only Memory), \textbf{PROM}
(Programmable ROM), \textbf{EPROM} (Erasable PROM), \textbf{EEPROM}
(Electrically Erasable PROM), and \textbf{flash memory}. Flash memory is the
modern choice because it can be erased and rewritten while still installed,
which is what makes firmware updates possible.

\begin{definitionbox}{Non-volatile memory}
Memory that retains its contents when the power is removed. ROM, flash
memory, an SSD, and an HDD are non-volatile. RAM is the common volatile
memory.
\end{definitionbox}

The important contrast is with RAM. RAM is volatile and is empty at
power-on; firmware is non-volatile and is present at power-on. A common wrong
statement is that firmware is loaded into RAM before the machine starts --- in
fact the firmware resides in ROM/flash and runs from there.

It is worth naming the members of the ROM family, because exam papers often
ask which memory type can be rewritten. \textbf{ROM} is programmed at the
factory and cannot be changed. \textbf{PROM} (Programmable ROM) can be
written once by the user. \textbf{EPROM} (Erasable PROM) can be erased with
ultraviolet light and reprogrammed. \textbf{EEPROM} (Electrically Erasable
PROM) can be erased and rewritten electrically. \textbf{Flash memory} is a
fast, block-erasable form of EEPROM and is the type used for modern firmware,
which is why firmware can be updated in place.

\begin{example}
An item states, ``The firmware of a modern PC is stored in RAM so that it can
be updated.'': This is wrong on the first point. Firmware is stored in
non-volatile ROM or flash memory, not in volatile RAM. Flash memory is what
makes the firmware updatable, not RAM.
\end{example}

\section{Hardware, software, and firmware}
It helps to place firmware among the three basic layers of a computer.

\begin{center}
\begin{tabular}{p{0.16\textwidth}p{0.36\textwidth}p{0.38\textwidth}}
\toprule
\textbf{Layer} & \textbf{What it is} & \textbf{Example} \\
\midrule
Hardware & The physical parts of the computer & Processor, motherboard, keyboard \\
Firmware & Software fixed in non-volatile memory that controls the hardware & BIOS, UEFI \\
Software & Programs that the user runs or that run on the OS & MS Word, a web browser \\
\bottomrule
\end{tabular}
\end{center}

Hardware is the body, software is the instruction set, and firmware is the
low-level bridge that tells the hardware how to begin. Firmware is software
by nature, but it is treated separately because it is stored on the device
and is not installed or removed like an ordinary program.

\section{Firmware on everyday devices}
System firmware is not unique to computers. Many devices carry firmware that
gives them their basic behaviour. A printer has firmware that controls
printing; a Wi-Fi router has firmware that runs its networking; a smartphone
has firmware plus an operating system; a modern television has firmware for
its tuner and menu. Updating such a device usually means updating its
firmware, again by flashing. This is why a firmware fix can repair a device
without any change to the user's applications.

\section{System firmware: the BIOS}
On a personal computer, the most familiar firmware is the \textbf{BIOS},
which stands for \textbf{Basic Input/Output System}. It is stored on a chip
on the motherboard. When the computer is switched on, the BIOS runs first. It
performs the power-on self-test, initialises the basic hardware, reads its
settings, and then loads the operating system.

The single most important exam fact on this topic is that the firmware which
initialises the hardware and loads the operating system is the BIOS. This is
directly tested (PYQ-28). The BIOS is sometimes confused with a device driver
or with the kernel; neither of these loads the operating system. A device
driver is loaded \emph{by} the operating system, and the kernel is part
\emph{of} the operating system.

\begin{definitionbox}{BIOS}
The Basic Input/Output System: the system firmware on a PC motherboard that
runs first at power-on, initialises the hardware, and loads the operating
system.
\end{definitionbox}

\section{The power-on boot sequence}
The machine follows a fixed chain each time it starts. Learning this chain
makes the roles of firmware, the OS, and drivers easy to separate.

\begin{enumerate}
  \item \textbf{Power on.} Electrical power reaches the motherboard.
  \item \textbf{Firmware runs.} The BIOS or UEFI begins executing from
    ROM/flash.
  \item \textbf{Hardware initialisation and POST.} The firmware tests and
    configures the basic hardware.
  \item \textbf{Operating system loads.} The firmware finds the boot device,
    reads the bootloader, and hands control to the OS.
  \item \textbf{Drivers load.} Once the OS is running, it loads the device
    drivers it needs.
  \item \textbf{Applications run.} User programs start on top of the OS.
\end{enumerate}

The chain can be written compactly as: power-on $\rightarrow$ firmware
$\rightarrow$ initialisation $\rightarrow$ OS $\rightarrow$ driver
$\rightarrow$ application. Any question that asks ``which comes first?'' is
answered by walking this chain from left to right.

The step in which the firmware hands control to the operating system is worth
a little more detail. The firmware locates the boot device according to its
configured boot order, reads a small program called the
\textbf{bootloader}, and runs it. The bootloader then loads the main part of
the operating system into memory. Under the legacy BIOS, the bootloader is
found through the MBR (Master Boot Record); under UEFI it is found through
the GPT (GUID Partition Table) scheme and a dedicated EFI system partition.

\begin{definitionbox}{Bootloader}
A small program that the firmware runs to load the operating system into
memory and start it.
\end{definitionbox}

\section{Boot devices and boot order}
The firmware does not decide by itself which device to start from. It follows
the \textbf{boot order} stored in its settings: a list such as hard disk
first, then optical drive, then USB. The firmware tries each device in turn
until it finds one that can start the operating system. The boot order is the
reason a machine can be made to start from a USB drive by moving that drive
to the top of the list, and it is one of the first things a technician checks
when a computer reports that no boot device is found.

\begin{example}
A computer is switched on and shows, ``No boot device found.'' Walking the
chain helps: the firmware runs, the POST passes, and the firmware then looks
through the boot order for a startable device. If the disk has failed or its
entry has moved down the list, no operating system is found and the firmware
stops with this message. The failing layer is the firmware's boot
configuration or the boot disk, not the application software.
\end{example}

\section{POST: the power-on self-test}
\textbf{POST} stands for \textbf{Power-On Self-Test}. It is one of the first
things the BIOS does after power-on. During POST the firmware checks that the
essential hardware --- CPU, RAM, keyboard, disk controller, and video --- is
present and responding. If a component fails, the machine may beep in a code,
show an error message, or refuse to continue.

POST is a step \emph{inside} the firmware's work. It runs \emph{before} the
operating system, not after it. A frequent trap is to place the POST after
the OS has started; the correct order is firmware (including POST) first,
then the OS.

\begin{definitionbox}{POST}
The Power-On Self-Test: the hardware check the BIOS runs at startup before
loading the operating system.
\end{definitionbox}

\section{BIOS setup and CMOS}
The BIOS keeps a small set of settings: the system date and time, the order
in which the machine should look for a boot device, hardware options, and
sometimes passwords. These settings are stored in a tiny battery-backed
memory chip called \textbf{CMOS} (Complementary
Metal-Oxide-Semiconductor). Because the chip is battery-backed, the clock
and boot order are remembered between sessions, even when the main power is
off.

The BIOS Setup utility (sometimes called the CMOS setup) is opened by
pressing a key during startup. That key is commonly \textbf{Del} or
\textbf{F2}, but the exact key varies by manufacturer. Inside setup, the
user can change the boot order, enable or disable built-in devices, and set
passwords.

\section{UEFI: the modern firmware}
\textbf{UEFI} stands for \textbf{Unified Extensible Firmware Interface}. It
is the modern replacement for the legacy BIOS and performs the same basic
job: it starts the hardware and loads the operating system. UEFI adds
several improvements.

\begin{itemize}
  \item It boots the machine \textbf{faster}.
  \item It supports very large disks through the \textbf{GPT} (GUID
    Partition Table) scheme, instead of the older MBR scheme with its
    practical limit of about 2 TB.
  \item It provides a \textbf{graphical} setup utility, often with mouse
    support.
  \item It offers \textbf{Secure Boot}, which helps prevent unauthorised
    software from loading at startup.
\end{itemize}

For exam questions, remember that ``the firmware that boots the OS'' may be
either BIOS or UEFI depending on the age of the system. Both initialise the
hardware and load the operating system; they are firmware, never drivers.

\section{BIOS versus UEFI}
The two system-firmware generations can be compared on a few clear points.

\begin{center}
\begin{tabular}{p{0.18\textwidth}p{0.34\textwidth}p{0.40\textwidth}}
\toprule
\textbf{Point} & \textbf{BIOS} & \textbf{UEFI} \\
\midrule
Full form & Basic Input/Output System & Unified Extensible Firmware Interface \\
Age & Legacy (older) & Modern \\
Boot speed & Slower & Faster \\
Disk support & MBR (about 2 TB limit) & GPT (very large disks) \\
Interface & Mostly text-based & Graphical \\
Security & No Secure Boot & Secure Boot available \\
\bottomrule
\end{tabular}
\end{center}

\section{Device drivers}
A \textbf{device driver} is software that lets the operating system
communicate with a specific hardware device. Each device --- a printer, a
graphics card, a network card, a sound card --- needs its own driver, because
a driver translates the general commands of the operating system into the
particular instructions that one device understands.

Drivers are loaded \textbf{by the operating system}, after the firmware has
started the machine. They are matched to the exact device and the exact OS
version. Because a driver is software, it is wrong to call it hardware, and
because it is loaded by the OS, it is wrong to say that a driver initialises
the hardware and loads the OS.

\begin{definitionbox}{Device driver}
Software that enables the operating system to communicate with a particular
hardware device by translating general OS commands into device-specific
instructions.
\end{definitionbox}

The three layers are easy to mix up, so keep their roles separate:

\begin{center}
\begin{tabular}{p{0.20\textwidth}p{0.30\textwidth}p{0.34\textwidth}}
\toprule
\textbf{Layer} & \textbf{Stored where} & \textbf{Role} \\
\midrule
Firmware (BIOS/UEFI) & On the hardware (ROM/flash) & Starts the machine; initialises hardware; loads the OS \\
Device driver & With the OS, on disk & Talks to one specific device \\
Operating system & On disk & Manages all resources and provides the user interface \\
\bottomrule
\end{tabular}
\end{center}

\begin{example}
A question asks, ``Which software initialises the hardware and loads the
operating system?'' A device driver is offered as an option. This is the
classic near-neighbour trap. The correct answer is the \textbf{BIOS/UEFI}
firmware. A device driver cannot do this job: it is itself loaded by the OS,
and it controls one device rather than starting the whole machine.
\end{example}

\section{Firmware, driver, OS, and kernel}
Four terms are often mixed up because they all describe ``software close to
the hardware''. Their roles are different, and the exam tests the
differences.

\begin{center}
\begin{tabular}{p{0.15\textwidth}p{0.40\textwidth}p{0.35\textwidth}}
\toprule
\textbf{Term} & \textbf{What it does} & \textbf{When it runs} \\
\midrule
Firmware (BIOS/UEFI) & Initialises hardware and loads the OS & First, at power-on \\
Kernel & The core of the OS; manages memory, processes, and devices & After the OS starts loading \\
Device driver & Communicates with one specific device & After the OS is running \\
Operating system & Manages all resources and provides the interface & After firmware, before applications \\
\bottomrule
\end{tabular}
\end{center}

The kernel is part of the operating system, not separate firmware. A driver
is loaded by the OS once the kernel is running. Only firmware runs before the
OS. If an item asks which layer ``starts the machine and loads the OS'', the
answer must be firmware --- BIOS or UEFI --- and never a driver or the
kernel.

\section{Firmware updates and flashing}
A \textbf{firmware update} replaces the firmware stored on a device with a
newer version. The process is commonly called \textbf{flashing}, because the
new firmware is written into flash memory. Updates are issued to fix bugs,
add features, and patch security flaws.

Firmware updates are much rarer than ordinary software updates, and they
carry more risk. Firmware is low-level, so a failed update can leave a device
unusable --- often described as ``bricking'' it. Interrupting the power
during flashing, or loading the wrong firmware for the model, can damage the
device permanently. For this reason the standard advice is to use only the
manufacturer's update for the exact model and to keep a stable power supply
throughout the update.

\begin{example}
A candidate reads, ``A firmware update is called formatting.'' This is wrong.
Formatting prepares a storage volume for files; it does not rewrite firmware.
The correct term for a firmware update is \textbf{flashing}.
\end{example}

\section{Exam retrieval and common confusions}
The following distinctions prevent most errors on this topic.

\begin{enumerate}
  \item Firmware that initialises hardware and loads the OS is
    \textbf{BIOS/UEFI}, not a device driver or the kernel (TRAP-24, PYQ-28).
  \item A \textbf{device driver} is \textbf{software}; the OS is system
    software (TRAP-10).
  \item Firmware is stored in \textbf{ROM/flash} and is non-volatile; RAM is
    volatile and empty at power-on.
  \item \textbf{POST} runs \textbf{before} the operating system, as part of
    the firmware's startup work.
  \item A driver is loaded \textbf{by the operating system}, after the
    firmware has started the machine.
  \item The startup order is power-on, firmware, initialisation, OS, driver,
    application.
  \item A firmware update is called \textbf{flashing}; interrupting it can
    damage the device.
\end{enumerate}

\begin{example}
An item asks which layer runs first when a computer is switched on. Trace the
chain: firmware runs before the OS, and the OS runs before any driver, which
runs before any application. The answer is the firmware (BIOS/UEFI). If the
stem instead asks what tests the hardware at startup, the answer is the POST;
if it asks what lets the OS talk to one device, the answer is the device
driver.
\end{example}

\subsection*{Exercises}
\begin{enumerate}[label=\arabic*.]
  \item Define firmware and state where it is normally stored.
  \item List the steps of the power-on boot sequence from power-on to
    application.
  \item Expand BIOS, UEFI, and POST, and give one role of each.
  \item Distinguish firmware, a device driver, and an operating system.
  \item What is a firmware update called, and what is the main risk during
    one?
\end{enumerate}

\subsection*{Solutions}
\begingroup\small
\begin{enumerate}[label=\arabic*.]
  \item Firmware is software stored in non-volatile memory that gives a
    device its basic, low-level startup and control instructions. It is
    normally stored in ROM or flash memory on the hardware itself.
  \item Power on $\rightarrow$ firmware (BIOS/UEFI) runs $\rightarrow$
    hardware initialisation and POST $\rightarrow$ the operating system
    loads $\rightarrow$ device drivers load $\rightarrow$ application
    programs run.
  \item BIOS = Basic Input/Output System: system firmware that initialises
    hardware and loads the OS. UEFI = Unified Extensible Firmware Interface:
    the modern firmware with fast boot, GPT disk support, and Secure Boot.
    POST = Power-On Self-Test: the hardware check the BIOS runs at startup.
  \item Firmware is stored on the hardware in ROM/flash and starts the
    machine and loads the OS. A device driver is software loaded by the OS
    that talks to one specific device. The operating system is the large
    system-software layer that manages all resources and provides the user
    interface.
  \item A firmware update is called flashing. The main risk is that
    interrupting the power during flashing, or using the wrong firmware, can
    leave the device unusable.
\end{enumerate}
\endgroup
\end{document}
